\documentclass[12pt]{article}
\usepackage[utf8]{inputenc}
\usepackage[T1]{fontenc}
\usepackage{hyperref}
\title{Yakut Language Resources for K-12 Curriculum Development}
\author{Compiled by Claude Sonnet 5.5 (claude-sonnet-5-5) for Andr\'as Kornai}
\date{}
\begin{document}
\maketitle

\section{Introduction}
Yakut (Sakha; ISO 639-3 sah; Glottocode yaku1245) is a Turkic language and an official language of the Sakha Republic (Yakutia) in the Russian Federation. It is taught as a school subject. The searches did not return a clear listing of current school textbooks, so this dossier is thin on primary materials and gives pointers instead.

\section{Primary Educational Resources}
\begin{itemize}
\item Open lesson plan for a Sakha language class from a village school (infourok.ru), an example of teacher-made classroom material: \url{https://infourok.ru/otkrytyj-urok-na-temu-da5aahyn-aat-5770828.html}
\item Public open letter from parents and community members about preserving and developing national-language schools in Sakha (a petition, not a teaching resource, but it shows the policy debate): \url{https://www.change.org/p/глава-республики-саха-якутия-николаев-а-с-открытое-обращение-главе-рс-я-николаеву-а-с-о-сохранении-и-развитии-национальных-школ}
\end{itemize}
Search results suggested that the Yakutsk publisher Bichik and the regional institute for education development (IRO RS(Ya)) are likely sources of textbooks, but no URLs for them were seen, so none is cited.

\section{Linguistic and Reference Materials}
\begin{itemize}
\item Universal Dependencies Yakut treebank (YKTDT): \url{https://universaldependencies.org/treebanks/sah_yktdt}
\item Tools for supporting language learning for Sakha (Apertium-based morphological analyzer, integrated into the Revita learning platform): \url{https://preview.aclanthology.org/naacl24-info/W19-6117}
\item World Loanword Database vocabulary for Sakha: \url{https://wold.clld.org/vocabulary/19}
\item Cambridge SPRI library catalogue record, apparently a Yakut-Russian dictionary of 1972: \url{https://www.spri.cam.ac.uk/library/catalogue/records/50510/}
\end{itemize}
The 15-volume explanatory dictionary of Yakut (completed 2018, about 80,000 entries) was mentioned in search summaries, but no direct URL for it was seen.

\section{Religious and Community Resources}
No scripture or community organisation resources were confirmed in these searches.

\section{Digital and Audio Resources}
\begin{itemize}
\item Self-published English-to-Yakut dictionary on Payhip (informational only, not authoritative): \url{https://payhip.com/b/5SBDF}
\end{itemize}
No audio or video resources were confirmed.

\section{Curriculum Development Context}
Yakut has official status and a presence in schools as a subject, a large speaker community and some computational resources. The main gaps found here are the identification of current textbook titles and a free online dictionary. A next step would be to check the Sakha Ministry of Education and Bichik publisher catalogues directly.

\end{document}